\documentclass[12pt]{article}
\usepackage[a4paper,margin=1in]{geometry}\usepackage[T1]{fontenc}
\usepackage{lmodern,microtype}\usepackage{amsmath,amssymb,amsthm,mathtools}
\usepackage{booktabs,array,enumitem,xcolor}\usepackage[numbers,sort&compress]{natbib}
\usepackage[colorlinks=true,linkcolor=blue!55!black,citecolor=blue!55!black,urlcolor=blue!55!black]{hyperref}
\linespread{1.07}
\newtheorem{theorem}{Theorem}[section]\newtheorem{proposition}[theorem]{Proposition}
\theoremstyle{remark}\newtheorem{remark}[theorem]{Remark}
\newcommand{\norm}[1]{\left\lVert#1\right\rVert}

\title{Gauge-Invariant Residue Cocycles and a Scalar Renormalization Obstruction\\
\large Physical Weights Along the Two Conjugate Branches}
\author{Liang Wang}\date{July 2026}
\begin{document}\maketitle
\begin{abstract}
The branch correspondence of RH-204 permits a gauge-free comparison of
source--observation residues across scales.  For matched nonzero residues
$r_{c,j},r_{f,j}$, the exact diagonal multiplier is
$m_j=r_{f,j}/r_{c,j}$.  These multipliers compose under refinement and occur
in conjugate pairs for real physical data.

We test whether the four multipliers can be replaced by one scale-dependent
complex scalar.  The least-squares optimum is explicit.  On the four
physical transitions its relative residual ranges from $0.32371$ to
$0.99985$; on the first left transition the optimal scalar is nearly zero
while the two branch multipliers have moduli $0.10961$ and $0.15002$ with
different phases.  The exact diagonal cocycle residual is below
$1.29\times10^{-16}$, and conjugation symmetry holds within
$3.23\times10^{-14}$.

Thus a common scalar normalization cannot transport the physical transfer
weights, even though a branch-diagonal endpoint cocycle always can.  The
cocycle is source/observation dependent and differs between left and right
channels by as much as $1.05364$; it is therefore not yet an intrinsic
spectral renormalization.  The negative result motivates separating the
weighted transfer ledger from the unweighted quartet divisor.
\end{abstract}

\section{Residues after branch matching}

For a simple eigenvalue $\lambda$ with right and left vectors $v,w$ chosen
so that $w^*v=1$, the spectral projector and physical residue are
\begin{equation}\label{eq:residue}
 P_\lambda=vw^*,\qquad
 r_\lambda=\operatorname{tr}(OP_\lambda S).
\end{equation}
RH-195 proves that this is also the Frobenius pairing of the source and
observation channel states \citep{WangRH195}.

RH-204 supplies the matched labels $j\in\{-,+\}$ for the two
upper-half-plane branches and their conjugates \citep{WangRH204}.  No
eigenvector phase is used in this labeling.

\section{Gauge invariance}

\begin{proposition}[Gauge invariance]\label{prop:gauge}
Under the biorthogonal rescaling
\begin{equation*}
 v\mapsto \alpha v,\qquad
 w\mapsto \overline\alpha^{-1}w,
 \qquad \alpha\ne0,
\end{equation*}
both $P_\lambda$ and $r_\lambda$ are unchanged.
\end{proposition}

\begin{proof}
The rescaled outer product is
$(\alpha v)(\overline\alpha^{-1}w)^*=vw^*$.  Substitute this identity into
\eqref{eq:residue}.
\end{proof}

Consequently, the large residue changes in RH-202 cannot be removed by
rephasing or balancing individual eigenvectors.

\section{Diagonal residue cocycle}

For every matched nonzero residue define
\begin{equation}\label{eq:multiplier}
 m_j(c,f)=\frac{r_{f,j}}{r_{c,j}}.
\end{equation}

\begin{theorem}[Finite cocycle law]\label{thm:cocycle}
For three levels $a,b,c$ on one consistently labeled branch,
\begin{equation}\label{eq:cocycle}
 m_j(a,c)=m_j(b,c)m_j(a,b).
\end{equation}
If the operators, source, and observation are real, then
\begin{equation}\label{eq:conjugate}
 m_{\overline j}(a,b)=\overline{m_j(a,b)}.
\end{equation}
\end{theorem}

\begin{proof}
Equation \eqref{eq:cocycle} is cancellation of the intermediate nonzero
residue.  Reality gives
$r_{\overline\lambda}=\overline{r_\lambda}$, and taking ratios gives
\eqref{eq:conjugate}.
\end{proof}

The diagonal cocycle is exact but endpoint dependent.  It does not predict a
new residue from coarse data alone.

\section{Best common scalar}

Let $r_c,r_f\in\mathbb C^4$ collect the ordered quartet residues.  Consider
\begin{equation}\label{eq:ls}
 \min_{\alpha\in\mathbb C}\norm{r_f-\alpha r_c}_2.
\end{equation}

\begin{proposition}[Scalar optimum]\label{prop:scalar}
If $r_c\ne0$, the unique minimizer is
\begin{equation}\label{eq:alpha}
 \alpha_*=\frac{r_c^*r_f}{r_c^*r_c}.
\end{equation}
The normalized residual vanishes if and only if $r_f$ is a scalar multiple
of $r_c$, equivalently all nonzero branch multipliers agree.
\end{proposition}

This is a strict, basis-free test of scalar scale renormalization.  A large
residual means that no global amplitude and phase correction can align all
physical weights.

\section{Physical multiplier table}

The upper-half-plane branch multipliers and scalar residuals are:
\begin{center}
\small
\begin{tabular}{llrrr}
\toprule
step & side & $m_-$ & $m_+$ & scalar residual\\
\midrule
$0.04\to0.02$ & left
& $0.1037+0.0356i$ & $-0.0269-0.1476i$ & $0.99985$\\
$0.04\to0.02$ & right
& $0.1009+0.1534i$ & $0.0372-0.1966i$ & $0.97654$\\
$0.02\to0.01$ & left
& $1.5204-0.3314i$ & $1.3920+0.6344i$ & $0.39406$\\
$0.02\to0.01$ & right
& $0.4963-0.0836i$ & $0.9778+0.3073i$ & $0.32371$\\
\bottomrule
\end{tabular}
\end{center}
The first transition is nearly orthogonal to a common-scalar description.
Even on the finer transition, one scalar leaves a $32$--$39\%$ residual.

The diagonal reconstruction
\begin{equation*}
 r_f=\operatorname{diag}(m_1,\ldots,m_4)r_c
\end{equation*}
has maximum relative residual $1.2846\times10^{-16}$, as required by the
exact definition.  Conjugate multiplier errors are below
$3.23\times10^{-14}$.

\section{Channel dependence}

At $0.04\to0.02$, the largest difference between corresponding left and
right branch multipliers is $0.11781$.  At $0.02\to0.01$ it is $1.05364$.
Therefore the multiplier is not currently universal across the two
physical realizations.

Across both transitions, the composed upper-branch multipliers are
\begin{center}
\begin{tabular}{lrr}
\toprule
side & $m_-(0.04,0.01)$ & $m_+(0.04,0.01)$\\
\midrule
left  & $0.16942+0.01977i$ & $0.05617-0.22251i$\\
right & $0.06290+0.06770i$ & $0.09681-0.18076i$\\
\bottomrule
\end{tabular}
\end{center}
The telescoping law is exact, but the two-step numbers reinforce rather than
remove branch and channel dependence.

This is expected from \eqref{eq:residue}: changing $S$ or $O$ changes the
weights without changing the eigenvalue divisor.  A transfer determinant
may need this cocycle, while an unweighted spectral determinant should not.

\section{Three distinct notions of renormalization}

The audit separates:
\begin{enumerate}[leftmargin=2em]
\item \emph{eigenvector gauge}, which leaves residues unchanged;
\item \emph{common scalar scale normalization}, rejected in the four cases;
\item \emph{branch-diagonal endpoint cocycle}, exact but source dependent.
\end{enumerate}
Conflating these notions would hide the obstruction by fitting one multiplier
per datum.

\section{Consequence for determinant design}

The weighted moments
\begin{equation}\label{eq:weighted}
 h_q=\sum_jr_j\lambda_j^q
\end{equation}
carry the residue cocycle.  The Newton traces
\begin{equation}\label{eq:unweighted}
 s_q=\sum_j\lambda_j^q
\end{equation}
do not.  RH-199 already proves that these two ledgers must not be identified.
The present obstruction strengthens the practical reason to study the
unweighted quartic divisor separately before attempting a physical feedback
factor.

\section{Claim boundary and next step}

The gauge and cocycle statements are exact finite algebra.  The rejection of
a common scalar is a four-case floating result, not an all-level theorem.
No residue lower bound, source-independent cocycle, von Mangoldt weight, or
arithmetic trace identity is obtained.  RH-207 next studies the unweighted
quartic divisor across channels and scales.

\section{What a useful residue theorem would require}

An all-level weighted transfer ledger would need more than nonzero endpoint
residues.  One needs a uniform lower bound preventing loss of observability,
an upper bound on cocycle products, and compatibility with the growing-cloud
selection.  If multipliers are allowed to be fitted independently at every
branch and level, exact reconstruction is tautological and provides no
compactness.  A useful theorem must derive them from source/observation
regularity and operator dynamics.

\bibliographystyle{plainnat}\bibliography{references}\end{document}
